\section{Ito Integral}
The solution depends on the specific form of the drift, $\alpha(X_s)$, and diffusion, $\sigma(X_s)$, the coefficients being linear or non-linear. The Fokker-Planck \cite{wiki:fokker-planck} partial differential equation (PDE) can be used to check if these coefficients are linear. Then apply either the integrating factor method or by using Ito's lemma.

We approximate a stochastic process $f(t)$ with a sequence of simple functions, $\{f_0(t),\cdots, f_n(t)\}$ where each  function is constant over the specific $\Delta t$ partitioned along $[0, T]$ time horizon. We define the time intervals as $(t_0 < t_1), (t_1< t_2), \cdots,(t_{n-1} < t_{n = T})$. As $n \rightarrow \infty$ the time interval , $\Delta t \rightarrow 0$.
With a simple function, we approximate the integral with respect to time interval as
\begin{align*}
\int_0^T f_n(t)dt \approx \sum_{k=0}^{n-1} f(t_k)\Delta t_k \qquad | \;\Delta t_k < (t_{k+1}-t_k) 
\end{align*}
The integral for the same function with respect to the Brownian process $W_t$
\begin{align*}
\int_0^T f_n(t)dW_t = \sum_{k=0}^{n-1} f(t_k)dW_t \int_{t_k}^{t_{k+1}} dW_t
\end{align*}
Using the property, $\int_{t_k}^{t_{k+1}} dW_t = \left(W_{t_{k+1}} - W_{t_k}\right) = \Delta W_k$ 
\begin{align*}
\int_0^T f_n(t)dW_t = \sum_{k=0}^{n-1} f(t_k)\Delta W_k
\end{align*}
The sequence of sums converges to a random variable, and for any small error $\epsilon > 0$ and $\delta > 0$, the probability that the difference between the sum and limit is greater than $\epsilon$ goes to zero as $n \rightarrow 0$. Similarly, the expected square difference between the sums and the limit goes to zero.
This is established using the Ito isometry;

\begin{align*}
\mathbb{E}\left[\left(\int_0^T dW_t \right)^2\right] = \mathbb{E}\left[\int_0^Tf(t)^2 dt\right]
\end{align*}

This property shows that the variance of the Ito integral \cite{wiki:ito-calculus} is equal to the expected value of the integral of the square of the integrand. It is instrumental in proving the convergence of the sum to the Ito integral in mean square.
\\\\
\subsubsection*{Example 1}
Consider the integral:
\begin{align}\label{eq:ito_integral_1}
I_t = \int_0^T f(t, W_t) dW_t
\end{align}

\begin{figure}[h!]
    \centering
    \includegraphics[width=0.5\textwidth]{./figures/ito_integral_1.pdf}
    \vspace{-25pt}
    \caption{\small Approximation of Ito integral}
    \label{fig:ito_integral_1}
\end{figure}
%\FloatBarrier

We approximate the solution using a discretized sum over a partition between $[0, T]$
\begin{align*}
I_t^{(n)} = \sum_{i=0}^{n-1} f(t_i, W_{t_i}) (W_{t_{i+1} - W_{t_i}})
\end{align*}
where the $dWt$ is (W_{t_{i+1} - W_{t_i}}) 
\sim \mathcal{N}(0, \Delta t)$.
\\\\
The solution to the integral using Ito's lemma is:
\begin{align*}
\int_0^T f(t, W_t) dW_t &= \int_0^T W_tdW_t\\
&= \frac{1}{2} W_t^2 - \frac{1}{2} T
\end{align*}
The numerical approximation used the Ito integral 
\begin{align*}
\int_0^T H_t dW_t &= L^2 - \lim_{n \rightarrow \infty}\int_0^T H_t (W_{t_{i+1}} - W_{t_i})
\end{align*}

where $H_t$ is adapted to the filtration generated by the Brownian process and that the integrand must be evaluated on the left endpoint of $t_i$. We simulate this solution and is shown in Fig. \eqref{fig:ito_integral_1}
\\\\
\subsubsection*{Example 2}
Consider an integral with a separate integrand $g(t)$

\begin{align}\label{eq:ito_integral_2}
\int_0^T g(t) dW_t
\end{align}

where $g(t)$ is a process defined by the SDE
\begin{align*}
dX_t = \mu dt + \sigma dW_t
\end{align*}

Let $\mu = 1$ and $\sigma=0.5$, this makes $dX_t=1dt + 0.5dW_t$ an arithmetic Brownian motion with drift.

The integrand $g(t)$ is a process of $X_t$, therefore, we will calculate $\int_0^T X_tdW_t$, using Ito's lemma  to $Y_t = f(X_t) = X_t^2$:

The solution to the integral using Itô's lemma is:
\begin{align*}
\int_0^T X_t \, dW_t = X_T^2 - X_0^2 - 2 \int_0^T X_t \, dt - 0.25 T
\end{align*}

This formula gives us a way to calculate the value of the Itô integral if we know 
$X_0$, $X_T$, and the integral of $X_t$ over time. We simulate the given SDE using the values $\mu=1.0$ and $\sigma=0.5$. We compare the exact solution using Ito's lemma versus the numerical simulated path as shown in Figure \eqref{fig:ito_integral_2}

\begin{figure}[h!]
    \centering
    \includegraphics[width=0.5\textwidth]{figures/ito_integral_2.pdf}
    \vspace{-25pt}
    \caption{\small Approximation of $\int_0^T g(t) dW_t$}
    \label{fig:ito_integral_2}
\end{figure}
%\FloatBarrier
